\section{Round-thirty-one reconstruction: resolvent-certified compression and admissible memory observations}
\label{sec:r31-a4}

The weak-Harris geometry remains on the exact complete-past state, but the
operator seminorm is weighted and local.  Orthogonal dynamics is now obtained
from a resolvent-small relative perturbation theorem rather than the bounded
perturbation theorem.  The resolved observation is Hilbert bounded, which
restores the advertised vertical-line remainder.  Finite-amplitude pressure
is continued only along certified bridges carrying a uniform spectral gap.

\subsection{Weighted history space and power drift}

Let $K(h,\dd m)$, $U(h,m)$, and $W$ be the A3 kernel, update, and Lyapunov
weight.  The fractional moment proved there gives
\[
 PW\le C W^\alpha,\qquad0<\alpha<1.                       \tag{A4.1}
\]
For $0<\sigma<1$ set
\[
 d_\sigma(h,h')=\sum_{j\ge0}\sigma^j
   (1\wedge d_E(m_{-j},m'_{-j})).                         \tag{A4.2}
\]
The transportation cost is capped,
\[
 d_H(h,h')=1\wedge
 \sqrt{d_\sigma(h,h')[1+W(h)+W(h')]},                     \tag{A4.3}
\]
but the operator space is not.  For $0<\gamma,\theta<1$ put
\[
 \norm f_{\cB_{\gamma,\theta}}
 =\sup_h{\abs{f(h)}\over1+W(h)^\gamma}
 +\sup_{h\ne h'}{\abs{f(h)-f(h')}\over
 d_\sigma(h,h')^\theta[1+W(h)^\gamma+W(h')^\gamma]}.      \tag{A4.4}
\]

\begin{theorem}[Weighted weak-Harris gap]
\label{thm:r31-a4-harris}
On the A3 certified family there are $C<\infty$, $\lambda<1$ such that
\[
 \norm{P^nf-\pi(f)}_{\cB_{\gamma,\theta}}
 \le C\lambda^n\norm f_{\cB_{\gamma,\theta}}.             \tag{A4.5}
\]
No finite-step transition law from one exact tail is required to dominate a
transition law from another exact tail.
\end{theorem}

\begin{proof}
Stable-holonomy distortion couples newly sampled suffixes and contracts
$d_\sigma$; the unchanged remote tails are discounted by the shift.  On a
Lyapunov sublevel a finite connector family gives a positive coupling
probability, while (A4.1) controls the complement.  Weak-Harris contraction in
(A4.3) follows.  Applying the same coupling to the weighted local quotient in
(A4.4), and using the power drift for distant pairs, yields the operator norm
estimate.  Since the denominator in (A4.4) grows with $W$, unbounded functions
of order $W^\gamma$ do not acquire the false bounded-oscillation property.
\end{proof}

\subsection{Small sources and certified finite-amplitude bridges}

For a source $g$ with
$|g(h)|\le g_0+\delta\log W(h)$ and the analogue of the seminorm in (A4.4),
choose $\alpha(\gamma+\delta)<\gamma$ and define
$P_gf=P(e^gf)$.

\begin{theorem}[Local fixed-space Feynman--Kac theorem]
\label{thm:r31-a4-local-fk}
On a neighbourhood of zero, $g\mapsto P_g$ is analytic on the single space
$\cB_{\gamma,\theta}$ and has a simple isolated positive leading eigenvalue
with analytic Riesz projection.
\end{theorem}

\begin{proof}
The pointwise estimate follows from
$P W^{\gamma+\delta}\le C W^{\alpha(\gamma+\delta)}$.
The weighted local quotient is controlled by splitting the multiplier and
kernel differences.  Normal convergence of
$\sum P(g^kf)/k!$ gives operator analyticity, and the isolated simple
projection at zero persists by analytic perturbation.
\end{proof}

A \emph{certified potential bridge} is a $C^1$ path
$t\mapsto g_t$, $0\le t\le1$, together with constants and a finite open cover
$I_1,\ldots,I_M$ such that on every $I_j$:

\begin{enumerate}[label=(B\arabic*)]
\item $P_{g_t}$ satisfies one common drift and weighted Lasota--Yorke
      inequality;
\item the predictive quotient has a uniform irreducibility/aperiodicity
      connector bound;
\item one contour $\Gamma_j$ separates a positive eigenvalue from the rest of
      the spectrum by a fixed distance.
\end{enumerate}

\begin{theorem}[Finite-amplitude continuation along a certified bridge]
\label{thm:r31-a4-bridge}
Along a certified bridge, the leading eigenvalue and Riesz projection remain
simple and $C^1$ (analytic when the bridge is analytic).  Every finite
exposing potential used in C2 is required to be connected to zero by such a
bridge.
\end{theorem}

\begin{proof}
The local theorem starts the continuation on $I_1$.  Conditions (B1)--(B3)
make the Riesz rank locally constant and positivity/irreducibility make it one.
On overlaps the projections agree by uniqueness.  A finite cover continues
the branch to $t=1$.  Large bounded multiplication alone is not used to infer
a spectral gap.
\end{proof}

\subsection{Typed suspension renewal}

The suspension state is
\[
 \Omega=\{(h,m,r):m\in\operatorname{supp}K(h,\cdot),
                       0\le r<\tau(m)\}.                  \tag{A4.6}
\]
Let $\cX$ be the weighted flight space and
$\cY=\cB_{\gamma,\theta}$.  For $\operatorname{Re}z$ large define
\begin{align}
 (R_0(z)f)(h,m,r)
 &=\int_0^{\tau(m)-r}e^{-zs}f(h,m,r+s)\dd s,              \tag{A4.7}\\
 (B(z)f)(h)
 &=\int K(h,\dd m)\int_0^{\tau(m)}e^{-zs}f(h,m,s)\dd s,  \tag{A4.8}\\
 (T(z)F)(h)
 &=\int K(h,\dd m)e^{-z\tau(m)}F(U(h,m)),                \tag{A4.9}\\
 (A(z)F)(h,m,r)
 &=e^{-z(\tau(m)-r)}F(U(h,m)).                            \tag{A4.10}
\end{align}

\begin{proposition}[Suspension resolvent identity]
\label{prop:r31-a4-renewal}
On the common resolvent set,
\[
 (z-L_{\rm sus})^{-1}
 =R_0(z)+A(z)(I-T(z))^{-1}B(z).                           \tag{A4.11}
\]
\end{proposition}

\begin{proof}
Integrate the resolvent equation along the current residual flight.  The
boundary value at its end is $A(z)$, each complete subsequent flight is
$T(z)$, and the forcing collected on a complete flight is $B(z)$.  The A3
roof moment justifies Fubini and the geometric series; analytic continuation
gives (A4.11).
\end{proof}

\subsection{Resolvent-small orthogonal compression}

Let $L_0$ generate a semigroup on $\cH$ satisfying
\[
 \norm{e^{tL_0}}_{\cH\to\cH}\le Me^{-\omega t}.            \tag{A4.12}
\]
Let $A:D(L_0)\to\cH$ be closed relative to $L_0$.  Assume that for some
$\sigma> -\omega$,
\[
 \sup_{\operatorname{Re}z\ge\sigma}
 \norm{A(z-L_0)^{-1}}\le\vartheta<1.                      \tag{A4.13}
\]

\begin{lemma}[Resolvent-certified relative perturbation]
\label{lem:r31-a4-relative}
The operator $L_0+A$ on $D(L_0)$ generates a semigroup, and
\[
 (z-L_0-A)^{-1}=(z-L_0)^{-1}
 [I-A(z-L_0)^{-1}]^{-1}.                                  \tag{A4.14}
\]
If (A4.13) holds on $\operatorname{Re}z\ge-\omega/2$, then
$\norm{e^{t(L_0+A)}}\le C e^{-\omega t/2}$.
\end{lemma}

\begin{proof}
The Neumann inverse in (A4.14) is uniform on the stated half-plane.  The
resolvent identity gives the Hille--Yosida powers by differentiating in $z$;
the inverse Laplace formula yields the semigroup and the exponential bound.
This is a relative-resolvent theorem, not the bounded perturbation theorem on
$\cH$.
\end{proof}

Let $P_0$ be a finite-rank spectral projection commuting with $L$, and let
$P$ be the physical resolved projection of equal rank.  The graph-angle map
$S=Q_0Q:Q D(L)\to Q_0D(L)$ is assumed invertible, and direct expansion gives
\[
 S(QLQ)S^{-1}=L|_{Q_0\cH}+A_P.                            \tag{A4.15}
\]
The finite-dimensional angle/coupling certificate is now the resolvent
condition
\[
 \sup_{\operatorname{Re}z\ge-\omega/2}
 \norm{A_P(z-L|_{Q_0\cH})^{-1}}<1.                        \tag{A4.16}
\]

\begin{theorem}[Stable compressed generator]
\label{thm:r31-a4-compression}
Under (A4.12), graph invertibility of $S$, and (A4.16),
$L_Q=QLQ$ generates an exponentially stable semigroup both in $\cH$ and on
its graph domain.
\end{theorem}

\begin{proof}
Apply Lemma~\ref{lem:r31-a4-relative} to (A4.15) and conjugate by $S$.
Graph-domain stability follows by applying the resolvent formula to $L_Qx$.
Every hypothesis is a finite-rank angle or resolvent estimate that can be
checked without pretending a graph-bounded perturbation is Hilbert bounded.
\end{proof}

\subsection{Admissible observations and memory expansion}

Let $\cR=P\cH$ and $B_Q=QLP$.  Choose the resolved basis
$e_1,\ldots,e_d$ so that $e_j\in D(L^*)$ and define
\[
 (C_Qx)_j=\ip{x}{Q L^*e_j},\qquad x\in\cH.                \tag{A4.17}
\]
Thus $C_Q\in\cL(\cH,\cR)$, even though its formal notation is $PLQ$.
Assume
\[
 B_Q\in\cL(\cR,D(L_Q^3)).                                 \tag{A4.18}
\]

\begin{theorem}[Admissible vertical-line memory calculus]
\label{thm:r31-a4-memory}
For $u(t)=Pe^{tL}x$,
\[
 u'(t)=PLPu(t)+\int_0^tK(t-s)u(s)\dd s+F_x(t),            \tag{A4.19}
\]
where
\[
 K(t)=C_Qe^{tL_Q}B_Q,\qquad F_x(t)=C_Qe^{tL_Q}Qx.          \tag{A4.20}
\]
Both decay exponentially.  On every vertical half-plane contained in the
compressed resolvent set,
\[
 \widehat K(z)=C_Q(z-L_Q)^{-1}B_Q                         \tag{A4.21}
\]
and
\[
 \widehat K(z)=z^{-1}C_QB_Q+z^{-2}C_QL_QB_Q
               +z^{-3}C_QL_Q^2B_Q+O(|z|^{-4}).           \tag{A4.22}
\]
The compressed resolvent satisfies the exact Schur complement
\[
 P(z-L)^{-1}P=
 [z-PLP-C_Q(z-L_Q)^{-1}B_Q]^{-1}.                         \tag{A4.23}
\]
\end{theorem}

\begin{proof}
Variation of constants in the $P/Q$ block equations gives
(A4.19)--(A4.21).  For $y=B_Qr\in D(L_Q^3)$, three resolvent identities give
\[
 (z-L_Q)^{-1}y=z^{-1}y+z^{-2}L_Qy+z^{-3}L_Q^2y
 +z^{-3}(z-L_Q)^{-1}L_Q^3y.
\]
Because $C_Q$ is bounded on $\cH$ and the stable resolvent is uniformly
Hilbert bounded, the last term is $O(|z|^{-3})$ before the additional retained
coefficient, and one further identity gives the $O(|z|^{-4})$ remainder in
(A4.22).  There is no graph-norm loss after applying $C_Q$.  Solving the
second resolvent block row proves (A4.23).
\end{proof}

Round Thirty's two operator counterexamples are therefore removed by explicit
resolvent-small relative generation and by a genuinely Hilbert-bounded
resolved observation.  Finite exposing potentials are used only when a
certified spectral bridge has been supplied.
